\section{Weekly Summary}

\begin{tcolorbox}[colback=yellow!10,colframe=black!50,title=AI-Generated Content]
\noindent The summary below was written by Claude (Anthropic) from that week's regression outputs and series descriptions. It has not been reviewed or fact-checked by a human, and should be read as an automated, informal recap -- not analysis.
\end{tcolorbox}

Welcome back to another thrilling week of shoving unrelated numbers into a blender and pressing "regress." As always, our machine roams the vast wilderness of FRED, grabs two series that have never met and have no business being in the same sentence, and demands to know how they feel about each other. Spoiler: mostly they don't, but a couple got weirdly clingy this week.

Monday kicked things off with a real crowd-pleaser: hospital and nursing care output versus house prices in Nassau-Suffolk, New York. The raw regression came out swinging with an R-squared of 0.60 and a p-value tiny enough to make a grad student weep tears of joy. "Long Island home values are driven by nursing homes!" screamed absolutely nobody who understands what a trend is. Because the moment we detrended it, the relationship folded like a lawn chair -- p-value of 0.23, effectively nothing. Classic trend mirage. Two things that both went up over the decades, holding hands only because time was passing. Cute, but no.

The midweek entries were mostly polite shrugs. Tuesday's government grants versus mining output deflator couldn't clear significance either way, which is honestly refreshing in its honesty. Then Wednesday gave us the week's genuine oddity: Canadian recession indicators versus the San Francisco Tech Pulse. This one actually survived detrending -- p-value still hilariously small at around 1.8e-11, with a negative relationship suggesting tech momentum slumps when Canada is officially in the dumps. Now, before anyone drafts a thesis, the R-squared is a humble 0.07, meaning the "relationship" explains roughly a rounding error's worth of anything. But it held up in both regressions, which around here counts as a standing ovation. My deeply rigorous theory: recessions are, shockingly, sometimes correlated with other recession-adjacent things. Nobel committee, my inbox is open.

Thursday's share-of-mortgages-held-by-the-top-1\% versus net government equipment depreciation was the plot twist of the week: barely significant raw, but the detrended version actually got stronger (p-value 0.005, R-squared jumping to 0.33) and flipped to a negative slope. Interesting enough to raise an eyebrow, meaningless enough to lower it again immediately. Friday's carbon emissions versus the discount rate did a similar small detrended flicker, while Saturday's equipment deflator versus Korean policy uncertainty pulled the reverse Monday -- looked great raw, evaporated once detrended. As ever, the standard disclaimer applies with full force: these are randomly paired series, chosen by a script with no opinions and no shame, and any "relationship" you see is spurious until proven otherwise, which we will never bother to do. Same nonsense next week.
